% The article of the trials of the text version of PDF (test/transcribing.cpp), set by pdfLaTeX
% three ways: with [T1]{fontenc} and lmodern (lmodern.pdf), with neither (ot1.pdf: OT1, cm
% fonts), and with [T1]{fontenc} alone (bitmap.pdf: EC fonts as bitmaps, cm-super missing). The
% preamble is \documentclass[twocolumn]{article} \usepackage[utf8]{inputenc} and those packages.
\section{Wstęp}
Pchnąć w tę łódź jeża lub ośm skrzyń fig. Źdźbło trawy, ćma i Łódź leżą na Śląsku, a żółw
płynie wolno. Zażółć gęślą jaźń\footnote{Pangram zawiera wszystkie polskie litery.}.
Efektywna finalizacja oficjalnych przepływów (ligatury fi, fl, ff, ffi).

\subsection{Lista}
\begin{itemize}
\item pierwszy punkt z wyrazem dzielonym na końcu wiersza: niekonstytucyjność,
\item drugi punkt, krótki,
\item trzeci punkt ze wzorem $E = mc^2$ i indeksem $x_{12}$.
\end{itemize}

\subsection{Tabela}
\begin{tabular}{lrr}
Miasto & Ludność & Rok \\
Gdańsk & 486 & 2024 \\
Olsztyn & 171 & 2024 \\
\end{tabular}

\section{Długi tekst}
Akapit 1. Wół gęśli źrebię, żubr chłonie ćmy; pięć łabędzi śpi nad źródłem, a sędzia ważył każdą głoskę osobno. Zdanie numer 1 kończy się kropką.

Akapit 2. Wół gęśli źrebię, żubr chłonie ćmy; pięć łabędzi śpi nad źródłem, a sędzia ważył każdą głoskę osobno. Zdanie numer 2 kończy się kropką.

Akapit 3. Wół gęśli źrebię, żubr chłonie ćmy; pięć łabędzi śpi nad źródłem, a sędzia ważył każdą głoskę osobno. Zdanie numer 3 kończy się kropką.

Akapit 4. Wół gęśli źrebię, żubr chłonie ćmy; pięć łabędzi śpi nad źródłem, a sędzia ważył każdą głoskę osobno. Zdanie numer 4 kończy się kropką.

Akapit 5. Wół gęśli źrebię, żubr chłonie ćmy; pięć łabędzi śpi nad źródłem, a sędzia ważył każdą głoskę osobno. Zdanie numer 5 kończy się kropką.

Akapit 6. Wół gęśli źrebię, żubr chłonie ćmy; pięć łabędzi śpi nad źródłem, a sędzia ważył każdą głoskę osobno. Zdanie numer 6 kończy się kropką.

Akapit 7. Wół gęśli źrebię, żubr chłonie ćmy; pięć łabędzi śpi nad źródłem, a sędzia ważył każdą głoskę osobno. Zdanie numer 7 kończy się kropką.

Akapit 8. Wół gęśli źrebię, żubr chłonie ćmy; pięć łabędzi śpi nad źródłem, a sędzia ważył każdą głoskę osobno. Zdanie numer 8 kończy się kropką.

Akapit 9. Wół gęśli źrebię, żubr chłonie ćmy; pięć łabędzi śpi nad źródłem, a sędzia ważył każdą głoskę osobno. Zdanie numer 9 kończy się kropką.

Akapit 10. Wół gęśli źrebię, żubr chłonie ćmy; pięć łabędzi śpi nad źródłem, a sędzia ważył każdą głoskę osobno. Zdanie numer 10 kończy się kropką.

Akapit 11. Wół gęśli źrebię, żubr chłonie ćmy; pięć łabędzi śpi nad źródłem, a sędzia ważył każdą głoskę osobno. Zdanie numer 11 kończy się kropką.

Akapit 12. Wół gęśli źrebię, żubr chłonie ćmy; pięć łabędzi śpi nad źródłem, a sędzia ważył każdą głoskę osobno. Zdanie numer 12 kończy się kropką.

